% K20
\section{System Design}\label{sec:system-design}

\subsection{Protocol}

The client holds a Paillier key pair; the server holds the public key and the model $(b, \mathbf{w})$ from \cref{eq:log-odds}.
Classifying one message takes a single round trip (\cref{tab:protocol}).
Instead of scoring the two classes separately, the server returns one encrypted log-odds value: this halves the server's work and the number of decryptions, gives the score used for ROC analysis directly, and reveals less of the model to the client.

\begin{table}[htbp]
  \centering
  \caption{One protocol run for a single message. $S$ is the fixed-point scale (\cref{sec:fixed-point}).}
  \label{tab:protocol}
  \begin{tabular}{@{}l l p{0.62\columnwidth}@{}}
    \toprule
    Step & Party & Operation \\
    \midrule
    1 & Client & Extract TF-IDF features $\mathbf{x} \in [0, 1]^k$ with the public vocabulary; encode $\tilde{x}_i = \lfloor S x_i \rceil$; send $c_i = \mathrm{Enc}(\tilde{x}_i)$ for \emph{all} $i = 1, \dots, k$. \\
    2 & Server & Encode $\tilde{w}_i = \lfloor S w_i \rceil$, $\tilde{b} = \lfloor S^2 b \rceil$; compute $c_\Delta = \mathrm{Enc}(\tilde{b}) \cdot \prod_{i} c_i^{\tilde{w}_i} \bmod n^2$ and return it. \\
    3 & Client & Decrypt $\tilde{\Delta} = \mathrm{Dec}(c_\Delta)$; label the message as spam if and only if $\tilde{\Delta} > 0$; use $\tilde{\Delta} / S^2$ as the score. \\
    \bottomrule
  \end{tabular}
\end{table}

By \cref{eq:hom-add,eq:hom-mul}, $c_\Delta$ is an encryption of $\tilde{\Delta} = \tilde{b} + \sum_i \tilde{w}_i \tilde{x}_i$, the log-odds computed on the encoded integers.
Negative weights need no special handling: \texttt{phe} represents them modulo $n$ and evaluates $c_i^{\tilde{w}_i}$ through a modular inverse.

\subsection{Fixed-Point Encoding}\label{sec:fixed-point}

Features and weights are scaled by $S = 10^{6}$ and rounded to integers; the bias is scaled by $S^2$, so that it has the same scale as the products $\tilde{w}_i \tilde{x}_i$.
We encode values ourselves instead of relying on the floating-point encoding of \texttt{phe} for two reasons.

\textit{No leakage through the exponent.}
For a float, \texttt{phe} chooses the public exponent attached to the ciphertext from the magnitude of the value, so it is not the same for every value: zero shares its exponent with values in $[1/16, 1)$, but smaller values and $1.0$ get different ones.
The server, which knows the encoding, would learn that the corresponding words occur in the message (\cref{sec:security}).
An integer always gets exponent~0, so all ciphertexts look alike.

\textit{Exact, verifiable computation.}
Over the integers, the homomorphic computation is exact: the decrypted $\tilde{\Delta}$ must be \emph{equal} to the same integer computation done in plaintext, which is our third success criterion (\cref{sec:threat-model}).
The only approximation is the rounding itself.
A zero feature is encoded exactly and contributes nothing, so the error of the score is bounded by a sum over the non-zero features only:
\begin{equation}
  \left| \frac{\tilde{\Delta}}{S^2} - \Delta \right| \le \frac{1}{2S} \sum_{i:\, x_i \ne 0} \left( |x_i| + |w_i| + \frac{1}{2S} \right) + \frac{1}{2S^2}.
\end{equation}
For SMS messages with a dozen non-zero features this is of the order of $10^{-5}$; we measure the actual error against the floating-point classifier in \cref{sec:results}.
The encoded values are also far from the plaintext limit: $|\tilde{\Delta}| \le S^2 (|b| + \sum_i |w_i|) < 10^{17}$, whereas \texttt{phe} accepts signed values up to $n/3 \approx 2^{2046}$ for a \num{2048}-bit key.

\subsection{Full Feature Vectors}

The client encrypts all $k$ features, including the zeros.
Sending only the non-zero positions would be cheaper but would reveal which vocabulary words occur in each message; we show in \cref{sec:security} that this alone is enough to classify spam.
The price is bandwidth: with \SI{512}{\byte} per ciphertext, a message costs \SI{256}{\kilo\byte} for $k = 500$ and \SI{512}{\kilo\byte} for $k = 1000$.

\subsection{Implementation}

The system is written in Python, and the model is trained with scikit-learn.
The weights $b$ and $\mathbf{w}$ are the differences between the spam and ham rows of \texttt{class\_log\_prior\_} and \texttt{feature\_log\_prob\_} of the fitted \texttt{MultinomialNB}, respectively.
The cryptographic operations use \texttt{phe} with \texttt{gmpy2}.
Client-side encryption ($k$ full encryptions per message) dominates the cost, so it runs in parallel across processes.
The server performs $k$ exponentiations with small exponents per message, and the client a single decryption.
